\section{Exercise 7: Radon-Nikodym's Derivative}

\textbf{Context:} Let $Q$ be a probability measure on $(\Omega, \mathcal{F})$ that is equivalent to $P$ and let $\mathcal{G} \subset \mathcal{F}$ be a sub-algebra.

\subsection{Part 1: Equivalence on the Sub-algebra}

\textbf{Question 1:} Show that $P$ \& $Q$ are equivalent on $(\Omega, \mathcal{G})$.

\textbf{Step-by-Step Explanation:}
To show that $P$ \& $Q$ are equivalent on the restricted space $(\Omega, \mathcal{G})$, we need to prove that for any event $A \in \mathcal{G}$, $P[A] = 0$ if and only if $Q[A] = 0$.

Because $\mathcal{G}$ is a sub-algebra of $\mathcal{F}$, any event $A$ that belongs to $\mathcal{G}$ automatically belongs to $\mathcal{F}$ ($A \in \mathcal{G} \implies A \in \mathcal{F}$).
We already know that $P$ \& $Q$ are equivalent on the larger space $(\Omega, \mathcal{F})$. This means that for any event in $\mathcal{F}$, they agree on what has a probability of zero.
Since every event in $\mathcal{G}$ is also in $\mathcal{F}$, the equivalence naturally carries over. Therefore, $P$ \& $Q$ are equivalent on $(\Omega, \mathcal{G})$.

Because of this equivalence, we can properly denote the Radon-Nikodym derivative defined specifically on $(\Omega, \mathcal{G})$ as $\frac{dQ}{dP}\big|_{\mathcal{G}}$.

\subsection{Part 2: Conditional Expectation Connection}

\textbf{Question 2:} Show that $\frac{dQ}{dP}\big|_{\mathcal{G}} = E^P\left[\frac{dQ}{dP} \middle| \mathcal{G}\right]$ P-a.s.

\textbf{Step-by-Step Explanation:}
To make the notation cleaner, let's define two variables:
Let $A = \frac{dQ}{dP}\big|_{\mathcal{G}}$ (the derivative restricted to $\mathcal{G}$).
Let $B = \frac{dQ}{dP}$ (the full derivative on $\mathcal{F}$).

We want to prove that $A = E^P[B | \mathcal{G}]$. 
By the mathematical definition of conditional expectation, $E^P[B | \mathcal{G}]$ is the unique $\mathcal{G}$-measurable random variable that satisfies the following property for every bounded $\mathcal{G}$-measurable test variable $Z$:
$$ E^P[B Z] = E^P\big[E^P[B | \mathcal{G}] Z\big] $$

Let's compute the expected value of $Z$ under the measure $Q$ in two different ways, using our two different Radon-Nikodym derivatives.

\textbf{Path 1 (Using B):} 
Because $Z$ is $\mathcal{G}$-measurable, it is also $\mathcal{F}$-measurable. Using the full derivative $B$ to change the measure from $P$ to $Q$ over $\mathcal{F}$, we get:
$$ E^Q[Z] = E^P[B Z] $$

\textbf{Path 2 (Using A):}
Because $Z$ is $\mathcal{G}$-measurable, we can use the restricted derivative $A$ to change the measure from $P$ to $Q$ directly over $\mathcal{G}$:
$$ E^Q[Z] = E^P[A Z] $$

\textbf{Conclusion:}
Since both paths compute the exact same expected value $E^Q[Z]$, we can equate them:
$$ E^P[A Z] = E^P[B Z] $$

By definition, $A$ is $\mathcal{G}$-measurable because it is the Radon-Nikodym derivative defined on the sub-algebra $\mathcal{G}$. 
Since $A$ is $\mathcal{G}$-measurable and satisfies $E^P[A Z] = E^P[B Z]$ for every bounded $\mathcal{G}$-measurable random variable $Z$, $A$ perfectly satisfies the definition of the conditional expectation of $B$ given $\mathcal{G}$. 

Due to the uniqueness of conditional expectation (up to a P-null set), we conclude:
$$ A = E^P[B | \mathcal{G}] \quad P\text{-a.s.} $$

Which translates exactly to our target equation:
$$ \frac{dQ}{dP}\bigg|_{\mathcal{G}} = E^P\left[\frac{dQ}{dP} \middle| \mathcal{G}\right] \quad P\text{-a.s.} $$